	\subsection{Cơ sở toán học của Sai số chuẩn vững (Robust Standard Errors)}
	\textbf{Khi nào cần dùng:} Khi các kiểm định trên đều cho ra $p\text{-value} < 0.05$ (Bác bỏ $H_0$, xác nhận có phương sai thay đổi). Lúc này, ta không vứt bỏ mô hình, mà dùng Robust SE để "cứu" các giá trị p-value của biến độc lập.
	
	\textbf{Cơ sở lý thuyết \& Chứng minh:} Xuất phát từ công thức hiệp phương sai thực tế của OLS:
	\begin{equation}
		\text{Var}(\hat{\boldsymbol{\beta}}) = (\mathbf{X}^{\top}\mathbf{X})^{-1} (\mathbf{X}^{\top} \boldsymbol{\Omega} \mathbf{X}) (\mathbf{X}^{\top}\mathbf{X})^{-1}
	\end{equation}
	Với $\boldsymbol{\Omega}$ là ma trận đường chéo chứa các $\sigma_i^2$ chưa biết. Huber \cite{huber1967} và White \cite{white1980} đã có một phát hiện toán học vĩ đại: Chúng ta không cần phải cố gắng ước lượng chính xác từng $\sigma_i^2$ riêng lẻ. Theo Luật số lớn (LLN), ma trận trung tâm $\frac{1}{n} \mathbf{X}^{\top} \boldsymbol{\Omega} \mathbf{X}$ có thể được ước lượng tiệm cận vững (asymptotically consistent) bằng cách thay thế trực tiếp phần dư mẫu vào: $\frac{1}{n} \sum \hat{\varepsilon}_i^2 \mathbf{x}_i \mathbf{x}_i^{\top}$.
	
	\textbf{Tối ưu hóa (HC0 đến HC3):}
	Cấu trúc kẹp Sandwich trên gọi là HC0. Tuy nhiên HC0 bị chệch trong mẫu nhỏ. MacKinnon và White (1985) \cite{mackinnon1985} đã tối ưu hóa nó thành phiên bản \textbf{HC3} bằng cách sử dụng các giá trị đòn bẩy $h_{ii}$ (thuộc ma trận hình chiếu Hat Matrix) để trừng phạt mạnh mẽ các quan sát ngoại lai (outliers):
	\begin{equation}
		\widehat{\text{Var}}_{HC3}(\hat{\boldsymbol{\beta}}) = (\mathbf{X}^{\top}\mathbf{X})^{-1} \left[\sum_{i=1}^{n} \frac{\hat{\varepsilon}_i^2}{(1-h_{ii})^2} \mathbf{x}_i\mathbf{x}_i^{\top}\right] (\mathbf{X}^{\top}\mathbf{X})^{-1}
	\end{equation}
	\textbf{Giải thích $h_{ii}$ -- Giá trị đòn bẩy (Leverage):}
	Đại lượng $h_{ii}$ là phần tử đường chéo thứ $i$ của \textbf{ma trận hình chiếu} (Hat Matrix), được định nghĩa:
	\begin{equation}
		\mathbf{H} = \mathbf{X}(\mathbf{X}^{\top}\mathbf{X})^{-1}\mathbf{X}^{\top}
	\end{equation}
	Ma trận này có tên gọi ``Hat Matrix'' vì nó ``đội mũ'' lên $\mathbf{y}$: giá trị dự đoán $\hat{\mathbf{y}} = \mathbf{H}\mathbf{y}$. Mỗi phần tử đường chéo $h_{ii}$ đo lường \textbf{mức độ ảnh hưởng (leverage)} của quan sát thứ $i$ trong không gian biến độc lập đến chính giá trị dự đoán $\hat{y}_i$ của nó. Cụ thể:
	\begin{itemize}
		\item $h_{ii}$ nằm trong khoảng $\left[\frac{1}{n},\; 1\right]$, với $\sum_{i=1}^n h_{ii} = p$ (số tham số của mô hình, bao gồm hệ số chặn).
		\item $h_{ii}$ lớn $\Rightarrow$ quan sát $i$ nằm xa trung tâm phân phối của $\mathbf{X}$ $\Rightarrow$ có khả năng ``kéo'' đường hồi quy về phía mình mạnh hơn.
		\item Quy tắc thực hành: quan sát có $h_{ii} > \frac{2p}{n}$ được coi là có đòn bẩy cao (high-leverage point).
	\end{itemize}
	Trong công thức HC3, hệ số $\frac{1}{(1-h_{ii})^2}$ đóng vai trò \textbf{trừng phạt} (penalize): những quan sát có leverage cao (nghĩa là $h_{ii}$ gần 1) sẽ khiến phần dư $\hat{\varepsilon}_i$ bị ``co lại'' so với sai số thực $\varepsilon_i$; việc chia cho $(1-h_{ii})^2$ bù đắp lại sự thiên lệch này, giúp ước lượng phương sai trở nên chính xác hơn trong mẫu nhỏ.

	HC3 được giới học thuật kinh tế lượng đánh giá là giải pháp \textbf{tối ưu nhất} và là tiêu chuẩn vàng khi xử lý phương sai thay đổi đối với dữ liệu cắt ngang (cross-sectional data).
